%-------------------------
% ATS-Friendly Resume Template
% Mimics Hong-Sheng Lai's resume style
% Compile with: pdflatex resume_template.tex
%-------------------------

\documentclass[letterpaper,11pt]{article}

\usepackage[utf8]{inputenc}
\usepackage[T1]{fontenc}
\usepackage{lmodern}
\usepackage[top=0.4in,bottom=0.4in,left=0.4in,right=0.4in]{geometry}
\usepackage{titlesec}
\usepackage{enumitem}
\usepackage{hyperref}
\usepackage{xcolor}

% ATS-friendly hyperlink setup
\hypersetup{
    colorlinks=true,
    linkcolor=black,
    urlcolor=blue!70!black,
    pdfauthor={Hong-Sheng Lai},
    pdftitle={Resume - Hong-Sheng Lai},
}

% Remove page numbers
\pagestyle{empty}

% Section formatting: bold title with horizontal rule
\titleformat{\section}
  {\large\bfseries}
  {}
  {0em}
  {}
  [\vspace{-0.75em}\rule{\textwidth}{0.4pt}]

\titlespacing*{\section}{0pt}{0.5em}{0.2em}

% Tighter list spacing
\setlist[itemize]{
  leftmargin=1.2em,
  itemsep=0pt,
  parsep=0pt,
  topsep=1pt,
  partopsep=0pt,
}

% No paragraph indent, ragged right to avoid word hyphenation
\setlength{\parindent}{0pt}
\raggedright

% --- Custom commands ---

% Entry header: {Left title}{Right location}
%               {Left subtitle}{Right date}
\newcommand{\entry}[4]{%
  \begin{tabular*}{\textwidth}{@{}l@{\extracolsep{\fill}}r@{}}
    \textbf{#1} & #2 \\
    #3 & #4 \\
  \end{tabular*}
}

% Project entry: {Title}{Venue/Affiliation}{Date}
\newcommand{\projectentry}[3]{%
  \begin{tabular*}{\textwidth}{@{}l@{\extracolsep{\fill}}r@{}}
    \textbf{#1} & \\
    #2 & #3 \\
  \end{tabular*}
}

% Paper entry: {Title}{Citation line (authors, venue, year, link)}
\newcommand{\paperentry}[2]{%
  \textbf{#1}\\[1pt]
  #2
}

%-------------------------
% DOCUMENT START
%-------------------------
\begin{document}

%----------NAME----------
\begin{center}
  {\LARGE\bfseries Hong-Sheng Lai}\\[2pt]
  +1(253)261-8683 \,$|$\,
  \href{https://github.com/hongsheng-lai}{github.com/hongsheng-lai} \,$|$\,
  \href{mailto:jl900613@gmail.com}{jl900613@gmail.com} \,$|$\,
  \href{https://www.linkedin.com/in/hongsheng-lai}{www.linkedin.com/in/hongsheng-lai}
\end{center}

%----------EDUCATION----------
\section{Education}

\entry
  {Johns Hopkins University}{Baltimore, MD}
  {PhD in Computer Science \textit{(advised by Dr.\ Steven Salzberg \& Dr.\ Mihaela Pertea)}}{Aug 2026 -- Present}
\entry
  {Carnegie Mellon University}{Pittsburgh, PA}
  {MS in Computational Biology \textit{(advised by Dr.\ Carl Kingsford)}}{Aug 2024 -- Dec 2025}
\entry
  {National Taiwan University}{Taipei, Taiwan}
  {BS in Biomechatronics Engineering, Minor in Economics \textit{(advised by Dr.\ Chien-Yu Chen)}}{Sep 2019 -- Jun 2023}

%----------SKILLS----------
\section{Skills}
\begin{itemize}
  \item \textbf{Programming Languages:} Python, C++, Java, SQL, R, LaTeX
  \item \textbf{MLOps \& Framework:} PyTorch, JAX, Grafana, Wandb
  \item \textbf{Developer Tools:} Git, Linux, Docker, Slurm, Kubernetes, CI/CD, MySQL, MongoDB, Apache Airflow
\end{itemize}

%----------EXPERIENCE----------
\section{Work Experience}

\entry
  {Genentech}{South San Francisco, CA}
  {Computational Biology Intern}{Jan 2026 -- Jul 2026}
\begin{itemize}
  \item Implemented a local \textbf{agentic AI} assistant using \textbf{LangGraph (ReAct)}, \textbf{LangChain}, and \textbf{Ollama} to let scientists configure sequencing analysis pipelines via natural language, featuring structured tool use with \textbf{Pydantic} schemas, real-time GUI control via a socket bridge, and on-premises inference for data privacy.
  \item Designed \textbf{Claude Code} skills for automating scientific document generation across 10 life-science domains, using structured prompt pipelines with multi-step citation verification (Crossref API), template-driven output, and built-in compliance checks.
  \item Developed \textbf{NGS} quality control pipelines for \textbf{cell therapy} materials in \textbf{long-read sequencing}, automating consensus generation, variant interpretation, and impurity inference, with results rendered in HTML/PDF templates.
\end{itemize}

\vspace{2pt}
\entry
  {Taiwan AI Labs}{Taipei, Taiwan}
  {Machine Learning Engineer Intern}{Jul 2023 -- Jan 2024}
\begin{itemize}
  \item Enabled sequence-function similarity search in \textbf{Qdrant vector database} with \textbf{contrastive learning} and \textbf{genomic language models}.
  \item Accelerated sparse matrix regression models with \textbf{JAX} for \textbf{federated analysis}.
  \item Established an end-to-end clinical annotation pipeline for \textbf{pharmacogenomics}-based drug usage recommendations.
\end{itemize}

\vspace{2pt}
\entry
  {Missmi Intelligence}{Taipei, Taiwan}
  {Software Engineer}{Jun 2021 -- Jun 2023}
\begin{itemize}
  \item Built and launched an NGS gene reporting product, collaborating with medical experts on action plan suggestions; over 400 units sold.
  \item Automated PDF report generation by integrating secondary analysis results into HTML templates, orchestrated with \textbf{Apache Airflow}.
\end{itemize}

\section{Projects}

\paperentry
  {ARCADE: Controllable Codon Design from Foundation Models via Activation Engineering}
  {\textit{ISMB -- MLCSB}, 2026. \href{https://www.biorxiv.org/content/10.1101/2025.08.19.668819v2}{[Paper]}}
\begin{itemize}
  \item Engineered a controllable codon design model by applying \textbf{activation engineering} techniques to both \textbf{decoder-only} and \textbf{encoder-decoder language models} to achieve \textbf{multi-objective optimization}.
\end{itemize}

\vspace{2pt}
\paperentry
  {Alignment-Based Copy Number Estimation and Allele Typing in Complex Genes}
  {\textit{In submission}. \href{https://www.biorxiv.org/content/10.1101/2025.09.13.675996v1}{[TypeAssembly]} \href{https://www.biorxiv.org/content/10.1101/2025.10.13.681303v2}{[deCYPher]}}
\begin{itemize}
  \item Developed \textbf{local alignment} algorithms for \textbf{copy number estimation} and \textbf{star-allele typing} in structurally complex gene regions (pharmacogenes, KIR, HLA) from haplotype-resolved \textbf{long-read assemblies}.
\end{itemize}

\end{document}
